\section{Finite stopping-band descriptors}\label{de:section}

Local erosion assigns a barrier to a neighborhood of each admissible
datum. We now choose finitely many of these neighborhoods covering all
but a prescribed amount of the reference probability. This is a finite
approximation in probability; the whole space of slit data need not be
compact, and no deterministic bound on every possible slit count is
required.

\begin{theorem}[Stopping-band transfer]\label{if:band-transfer}
Fix an inward or outward band, target, and ambient square chart as in
Section~\ref{if:section}, with all the stated positive clearances.
Use the critical plus-boundary law in the fixed larger reference square.
For every exceptional probability
$\epsilon>0$ and target error $\eta>0$, generic contour levels can be
chosen in any prescribed separated subwindows so that the following
hold for all sufficiently fine meshes.
\begin{enumerate}
\item The stopping observation of Lemma~\ref{if:stopping} has a finite
      descriptor $Z^\delta\in\{0,1,\ldots,m\}$, computed from band
      spins.  State $0$ is exceptional and has probability at most
      $\epsilon$ in the fixed reference chart.
\item Every good state $i$ has a deterministic regular signed barrier
      $\Gamma_i$.  For every realizable transcript with $Z^\delta=i$,
      adding its plus pins and then its minus pins changes the target
      law by a sum of at most $\eta$ in total variation.  The resulting
      target-side Gibbs law, denoted $Q_{i,\delta}$, is the same for
      every such transcript.  This assertion is uniform over all
      bounded observables on the target.
\item In the inward case one may also move the full barrier to a
      slightly farther signed contour, retaining the entire enclosed
      free component as the kept volume.  Counting the two old-pin
      removals as well, the sum of target errors can still be made
      at most $\eta$.  The final kept volume has a positive
      macroscopic buffer from its fixed boundary.  All geometric
      clearances have positive minima over the good states.
\item For a fixed finite list of bands the descriptors have convergent
      joint laws.  The bins can be chosen to be continuity sets of
      the limiting data law; hence this finite-space convergence is
      in total variation.
\end{enumerate}
The number of states, their barriers, and their positive clearances
depend on the fixed geometry and on $\epsilon,\eta$, not on the mesh.
The threshold for the mesh is chosen afterward.  If a different law,
possibly conditional on remote pins, has band density at most $K$
times that of the reference chart by
Corollary~\ref{fd:buffered-ratio}, its exceptional probability is at most
$K\epsilon$.  The same applies to mixtures of those conditional laws.
\end{theorem}

\begin{proof}
Choose generic levels by Lemma~\ref{if:visible-data}.  At each admissible
limiting datum apply Lemma~\ref{if:erosion}, and, for the inward version,
Lemma~\ref{if:barrier-buffer}.  Allocate the requested total-variation
budget among the finitely many ordered steps.  This produces an open
neighborhood and a barrier construction valid uniformly for every
realizable datum in that neighborhood on all sufficiently fine meshes.

Refine this cover to balls whose closures lie in such neighborhoods
and whose boundaries have limiting probability zero.  For a fixed
center, only countably many radii can give positive boundary mass.
Separability from Remark~\ref{if:metric-separable} supplies a countable
such subcover. By continuity of probability along its increasing finite
unions, take finitely many balls whose union has limiting
probability larger than $1-\epsilon/2$.  Order them and form disjoint
bins by assigning a datum to the first ball containing it.  Their
boundaries are contained in the finite union of the selected sphere
boundaries and therefore have probability zero.  The remaining bin
is exceptional.  Convergence of the data makes its lattice probability
at most $\epsilon$ for all sufficiently fine meshes.

Associate to a good bin the barrier of the neighborhood containing
its ball.  The local lemmas hold for every transcript in that bin,
with no dependence on an observable.  Take the minimum of the
finitely many positive clearances and mesh thresholds.  This proves
the first three assertions.  For a finite list use the same argument
in each coordinate and Lemma~\ref{if:visible-data} jointly.  Every
descriptor string is a continuity set, proving convergence of all
its finitely many probabilities.  Finally the bad event depends
only on band spins, so the density comparison gives the last
assertion directly, including after conditioning on the stated
remote information.
\end{proof}

\begin{corollary}[Lookup approximation of conditional averages]
\label{if:descriptor-averages}
In the outward setting of Theorem~\ref{if:band-transfer}, let
$\mathcal F^\delta$ be the stopping information.  For any family
of target observables $h_\delta$ satisfying $|h_\delta|\le H$, define
\[
       a_i(\delta)=\mathbb E_{Q_{i,\delta}}h_\delta
             \quad (i\ge1),
\]
and choose any exceptional-bin value in $[-H,H]$.  Then, on every
good transcript,
\begin{equation}\label{if:lookup-pointwise}
  \left|\mathbb E[h_\delta\mid\mathcal F^\delta]
             -a_{Z^\delta}(\delta)\right|\le 2H\eta.
\end{equation}
In particular, in the reference chart,
\begin{equation}\label{if:lookup-ltwo}
  \mathbb E\left|\mathbb E[h_\delta\mid\mathcal F^\delta]
             -a_{Z^\delta}(\delta)\right|^2
                  \le 4H^2(\eta^2+\epsilon).
\end{equation}
The good bins, barriers, and mesh threshold work for all these
observables simultaneously.  Neither the observables nor their
lookup values need converge as $\delta\downarrow0$.  More explicitly,
fix a reference continuation law $\nu$ and regard
$f_\delta^\nu(T)=\mathbb E_\nu[h_\delta\mid\mathcal F^\delta=T]$
as a function of the stopping transcript $T$.  Evaluating this same
function under any buffered band distribution with density bound
$K$ gives the integrated estimate with exceptional term $K\epsilon$.
It is the reference conditional-average function that is transferred.
It also equals the conditional average for the new law when that
law has the same continuation and distant wall given $T$.
If the distant wall is changed to an assignment $b$, the uniform
wall extension in Lemma~\ref{if:erosion} supplies the corresponding
completed laws $Q_{i,\delta}^{b}$. The conditional averages for this
continuation use the separate lookup values
$a_i^{b}(\delta)=\mathbb E_{Q_{i,\delta}^{b}}h_\delta$.
\end{corollary}

\begin{proof}
The conditional target law of a stopping transcript is its Gibbs
continuation by Lemma~\ref{if:stopping}.  On a good bin its total
variation distance from the single law $Q_{i,\delta}$ is at most
$\eta$.  The defining inequality for total variation gives
Equation~\eqref{if:lookup-pointwise}.  On the exceptional bin both
quantities are bounded by $H$ in absolute value.  Squaring and
integrating proves Equation~\eqref{if:lookup-ltwo}.  All constructions
in Theorem~\ref{if:band-transfer} preceded the choice of $h_\delta$.
The pointwise good-bin estimate and the transferred exceptional
probability prove the last assertion.  This is the reason that
mesh-dependent likelihoods can be transferred without assigning
them a continuum limit.
\end{proof}

\begin{corollary}[Random levels and finite ball probes]
\label{if:random-levels}
Theorem~\ref{if:band-transfer} and
Corollary~\ref{if:descriptor-averages} remain valid when the two contour radii are
drawn with a bounded density in a fixed compact rectangle of separated
radii lying strictly inside the permitted windows.  Different band
observations may use independent draws.  Their joint data laws
converge, including the auxiliary randomness.

Fix a countable dense collection of marked rational-walk prototypes
in the ambient space $\widehat{\mathcal X}$ of
Definition~\ref{if:datum}; probe balls are intersected with the
admissible-data space $\mathcal X_{\mathrm{band}}$. Allow a probe to test
whether the data lie within distance $r$ of one prototype, where
$r$ is drawn uniformly in a rational interval with closure in $(0,1)$.
For any prescribed $\epsilon,\eta$, finitely many such probes suffice
to obtain the lookup approximation in
Corollary~\ref{if:descriptor-averages}, with adjusted exceptional
probability at most $\epsilon$.  The finite probe list depends only
on the reference local geometry, the radius rectangle, and the
requested tolerances.  It is independent of the bounded observables
and of the size of an ambient buffered array.
\end{corollary}

\begin{proof}
For almost every radius pair the deterministic continuity argument
in Lemma~\ref{if:visible-data} applies.  Couple the radius draws
identically on all meshes, independently of the converging loops.
The records take values in the union over the permitted radii specified
in Definition~\ref{if:datum}. The deterministic rounding rule is
measurable in those draws, and its recording maps $R_\delta$ have
uniformly vanishing displacement. Thus the continuous mapping theorem,
or dominated convergence after this coupling, proves joint convergence
of the encoded data and the radius draws. The full stopping information
includes those draws. The compact radius rectangle
provides a uniform positive separation of contours and uniform
clearance from the window edges.  Thus the local erosion construction
still has a neighborhood for every admissible resulting datum.

For each datum choose a neighborhood $U_d$ as in
Lemma~\ref{if:erosion}.  By density, there are a prototype $p$ and
rational numbers $0<a<b<1$ such that
\[
                  d\in B(p,a),\qquad B(p,b)\subset U_d.
\]
Choose a rational threshold interval strictly between $a$ and $b$.
A positive test then certifies membership in $U_d$, while every
datum in $B(p,a)$ tests positive for every threshold in that interval.
The inner balls form an open cover of the admissible data.  Select
finitely many covering probability at least $1-\epsilon/2$ in the
reference law, and decode by taking the first positive test in
their fixed order.  If no test is positive, use the exceptional
state.  For a positive test use the single barrier law belonging
to its neighborhood $U_d$; hence
Equation~\eqref{if:lookup-pointwise} holds without dependence on
the random threshold or on $h_\delta$.

Every probe tie has probability zero because its threshold is an
independent continuous random variable, conditional on the data.
The finite string of bits consequently converges in law, and the
exceptional probability is at most $\epsilon$ for fine meshes.
Only finitely many entries in any fixed dense enumeration of the
prototypes and rational intervals have been used.  For finitely
many permissible radius windows, take the union of the corresponding
lists.  Finally the choices were made in one fixed local reference
chart.  Density comparison transfers their error bounds to other
buffered chart laws with the same fixed multiplicative constant,
regardless of the number of tiles in the larger array.
\end{proof}
